\section{The Scaling Hypothesis}

Discussing the realations between critical exponents is a foundamental step in the study and comprehension of critical phenomena. The scaling hypothesis is a powerful tool that allows us to derive these relations and understand the behavior of physical quantities near the critical point.

\dots

\[
    f(t,h) = \text{min}_m \left[ \frac{t}{2}m^2 + u m^4 - h \cdot m \right] = \begin{dcases}
        ... \\
        ...
    \end{dcases}
\]

\dots

Nom taking a generalized homogeneous function of two variables \(f(x,y)\) such that:
\[
    f(\lambda^{d_1} x, \lambda^{d_2} y) = \lambda f(x,y),
\]
if we set
\[
    \lambda = \left(\frac{1}{x}\right)^{\frac{1}{d_1}},
\]
we can notice how the function \(f\) can be rewritten as
\[
    f(1, \frac{1}{x^{d_2/d_1}} y) = \left(\frac{1}{x}\right)^{\frac{1}{d_1}} f(x,y).
\]
Now if we define \(g(z)=f(1,z)\), we can invert the relation above to obtain
\[
    f(x,y) = x^{d_1} g\left(\frac{y}{x^{d_2/d_1}}\right),
\]
which therefore shows how the function \(f\) can be expressed in terms of a single variable, which is the ratio \(\frac{y}{x^{\Delta}}\) (where \(\Delta = \frac{d_2}{d_1}\)) of the two original variables, up to an overall prefactor \(x^{d_1}\). This is the essence of the scaling hypothesis, which states that near the critical point, physical quantities can be expressed as homogeneous functions of the relevant variables, leading to scaling relations between critical exponents; we will see this more explicitly in the following sections.

If we assume that the true free energy density \(f\) is a generalized homogeneous function of the reduced temperature \(\vert t \vert\) and the external field \(h\) when we study it close to the critical point \(T_c\), we can write
\[
    f(\vert t \vert, h) = \vert t \vert^{d_1} g\left(\frac{h}{\vert t \vert^{\Delta}}\right),
\]
for some \(d_1\), \(\Delta\) and \(g\) (notice how we are now treaing \(d_1\) and \(\Delta\) as indipendent parameters). Knowing the singularities of the free energy density \(f(t,h)\), we can fix the values of \(d_1\) and \(\Delta\) and therefore derive the scaling relations between critical exponents.

For \(h=0\), we must have
\[
    f(t,0) \sim \frac{\vert t \vert^2}{u}, \quad \implies \quad d_1 = 2, \quad \lim_{x \to 0} g(x) \sim \frac{1}{u}.
\]
Furthermore, for \(t \to 0\), we must have
\[
    \lim_{t \to 0} f(t,h) \sim \frac{h^{4/3}}{u^{1/3}},
\]
which can be verified by setting
\[
    \lim_{x \to \infty} g(x) \sim \frac{x^{4/3}}{u^{1/3}}.
\]
Now, using the last relation, we can fix the value of \(\Delta\) as follows:
\[
    \lim_{t \to 0} f(t,h) \sim \vert t \vert^2 \frac{h^{4/3}}{u^{1/3} \, t^{4\Delta/3}}, \quad \implies \quad \Delta = \frac{3}{2},
\]
concluding that the homogeneity assumption in the free energy density \(f\) is consistent with the singularities predicted by the saddle point approximation, with
\[
    d_1 = 2, \quad \Delta = \frac{3}{2}.
\]
Next we are going to make the same assumptions for the true free energy density \(f\) and see how we can derive the scaling relations between critical exponents from it.\TODO{there were some errors in the notes, check them and fix them}

We assume that the true singular part of the free energy density \(f\) is a generalized homogeneous function of the reduced temperature \(\vert t \vert\) and the external field \(h\) when we study it close to the critical point \(T_c\), we can write
\[
    f(\vert t \vert, h) = \vert t \vert^{d} g\left(\frac{h}{\vert t \vert^{\Delta}}\right),
\]
for some \(d\), \(\Delta\) and \(g\). This does not imply the free energy of the mean field theory, which is a particular approximation of the true free energy (which can even have this form in general, but with different exponents), but it is a more general assumption that we make on the true free energy itself. Let us now find a way to relate the exponents \(d\) and \(\Delta\) to the critical exponents.

We start by relating \(d\) with the critical exponent \(\alpha\) that describes the singularity of the heat capacity \(C\) at the critical point. We know that
\[
    C = -T \frac{\partial^2 f}{\partial T^2},
\]
so that we have to compute the second derivative of the free energy density \(f\) with respect to the temperature \(T\). Computing the first derivative of \(f\) with respect to \(T\):
\[
    \frac{\partial f}{\partial T} \sim (2-\alpha) \dots
\]
\[
    = \, : \vert t \vert^{1-\alpha} \tilde{g}\left(\frac{h}{\vert t \vert^{\Delta}}\right),
\]
since in the previous parrenthesis \(h\) appeared always divided by \(\vert t \vert^{\Delta}\).

\begin{remark}
    More generally, partial derivative of an homogeneus generalized function will again be an homogeneus generalized function, with different prefactors and a different function \(\tilde{g}\) (which is related to the original function \(g\) by differentiation).
\end{remark}

Repeating the same procedure to compute the second derivative of \(f\) with respect to \(T\), we obtain
\[
    C \propto - \frac{\partial^2}{\partial t^2} f \sim \vert t \vert^{-\alpha} \tilde{\tilde{g}}\left(\frac{h}{\vert t \vert^{\Delta}}\right),
\]
where \(\tilde{\tilde{g}}\) is another appropritely defined function. Therefore, we can conclude that the critical exponent \(\alpha\) is related to the exponent \(d\) by the following relation:
\[
    \lim_{x \to 0} \tilde{\tilde{g}}(x) = \text{const} \implies \alpha = 2 - d,
\]
and thus we found the relation between the critical exponent \(\alpha\) and the exponent \(d\) that appears in the homogeneity assumption of the free energy density \(f\): \(d = 2 - \alpha\). The logic behind this is that we need to now how this function behaves at the critical point, where it could go to 0, diverge to infinity or go to a constant; since
\[
    \lim_{t \to 0} C (t, h=0) \sim \vert t \vert^{-\alpha},
\]
then we need to have \(\lim_{x \to 0} \tilde{\tilde{g}}(x) = \text{const}\) in order to have the correct singularity of the heat capacity \(C\) at the critical point.

Now the only unknown parameter is \(\Delta\), which we call the \textbf{gap exponent}: it is related to the critical exponent \(\beta\) that describes the singularity of the order parameter \(m\) at the critical point. We are saying that it is not important the theory in which we compute the critical exponent \(\beta\) (mean field, Landau-Ginzburg, etc), but it is important that the true free energy density \(f\) has the homogeneity property that we assumed above, with some value of \(\Delta\). Thus let us relate \(\Delta\) to \(\beta\) by computing the order parameter \(m\) as a function of the reduced temperature \(t\) at zero external field \(h=0\): remember that
\[
    m(t,h=0) \sim \vert t \vert^{\beta},
\]
then we can compute \(m\) by taking the derivative of the free energy density \(f\) with respect to the external field \(h\):
\[
    m = - \frac{\partial f}{\partial h} \sim \vert t \vert^{d-\Delta} g^{\prime} \left(\frac{h}{\vert t \vert^{\Delta}}\right),
\]
and therefore we can conclude that the critical exponent \(\beta\) is related to the gap exponent \(\Delta\) by the following relation:
\[
    \lim_{x \to 0} g^{\prime}(x) = \text{const} \implies \beta = d - \Delta,
\]
which is the second scaling relation that we were looking for
\[
    \beta = d - \Delta = 2 - \alpha - \Delta.
\]

now we can consider the critical exponent \(\delta\), which describes the singularity of the order parameter \(m\) at the critical point as a function of the external field \(h\):
\[
    m(t=0,h) \sim h^{1/\delta},
\]
then we can assume that the function \(g\) behaves as
\[
    \lim_{x \to \infty} g(x) \sim x^{p+1}, \quad \implies \quad \lim_{x \to \infty} g^{\prime}(x) \sim x^p,
\]
which is the only reasonable behavior that we can assume for the function \(g\) in order to have a power law singularity of the order parameter \(m\) at the critical point (indeed we need \(\lim_{x \to \infty} g^{\prime}(x) = \text{const}\)). Therefore, we can compute the order parameter \(m\) at the critical point \(t=0\) as a function of the external field \(h\):
\[
    \lim_{t \to 0} m(t,h) \sim \vert t \vert^{2-\alpha-\delta} \left(\frac{h}{\vert p \vert^{\Delta}}\right)^{p},
\]
which is well-defined only if \(p \Delta = 2 - \alpha - \delta\) leading us to
\[
    \lim_{t \to 0} m(t,h) \sim h^{\frac{2-\alpha-\delta}{\Delta}},
\]
which recalling the definition of the critical exponent \(\delta\) leads us to the following scaling relation:
\[
    \delta = \frac{\Delta}{2-\alpha-\Delta} = \frac{\Delta}{\beta}.
\]

For the susceptibility \(\chi\), we can use the same arguments and procedures, computing it as the derivative of the magnetization \(m\) with respect to the external field \(h\):
\[
    \chi(t,h) \propto \frac{\partial m}{\partial h} \sim \vert t \vert^{2-\alpha-2\Delta} g^{\prime}\left(\frac{h}{\vert t \vert^{\Delta}}\right),
\]

[...]

obtaining
\[
    \gamma = 2\Delta + \alpha - 2.
\]


After these computations, we have found the following scaling relations between critical exponents:
\[
    \begin{dcases}
        \alpha = 2 - d,                          \\
        \beta = d - \Delta,                      \\
        \delta = \frac{\Delta}{2-\alpha-\Delta}, \\
        \gamma = 2\Delta + \alpha - 2,
    \end{dcases}
\]
showing how the gap exponent \(\Delta\) appear in all the expressions for the critical exponents; now from these relations we can derive the following identities:
\[
    \begin{dcases}
        \alpha + 2\beta + \gamma = 2, \text{ Rushbrooke's identity,} \\
        \delta - 1 = \frac{\gamma}{\beta}, \text{ Widom's identity,} \\
    \end{dcases}
\]
[...]
\[
    \delta - 1 = \frac{\Delta}{2 - \alpha - \Delta} - 1,
\]
\[
    \frac{\gamma}{\beta} = \frac{\gamma \delta}{\Delta} = \frac{2\Delta + \alpha - 2}{\Delta} \left(\frac{\Delta}{2 - \alpha - \Delta}\right),
\]
[...]

These relations can be tested experimentally and numerically, and they are indeed satisfied with good precision, despite the fact that the individual exponents are different than the one predicted by the mean field theory.

    [\textit{values of exponents}]

So it seems that even in systems with very different exponents, the scaling relations are still satisfied, which is a strong evidence in favor of the scaling hypothesis and the homogeneity assumption that we made on the true free energy density \(f\).

\subsection{Correlation Functions}

The homogeneity assumption do not apply only to the free energy density \(f\), but also to other physical quantities derived from it. But it says nothing about the correlation functions. In order to obtain identities related to exponents of the correlation functions, we need to make two additional assumptions which will seem to be more foundamental and more general than the homogeneity assumption on the free energy density \(f\):
\begin{itemize}
    \item The correlation length \(\xi\) is a generalized homogeneous function of the reduced temperature \(\vert t \vert\) and the external field \(h\) when we study it close to the critical point \(T_c\):
          \[
              \xi(t,h) = \vert t \vert^{-\nu} g\left(\frac{h}{\vert t \vert^{\Delta}}\right).
          \]
    \item Close to criticality, \(\xi\) is the most important lenght scale in the system and is solely responsible for the singular contributions to thermodynamic quantities.
\end{itemize}

[...]

\subsection{Anomalous Dimension}

[...]